\chapter{Результаты}\label{ch:ch5}

\section{Сравнение алгоритмов построения линейных порядков}
\label{sec:ch5/indep-dep}

Все представленные алгоритмы обхода — \ref{alg:topo},
\ref{alg:topo-indep}, \ref{alg:topo-indep-lookup} и
\ref{alg:topo-dep} -- имеют одинаковую асимптотическую сложность в
общем случае. Однако алгоритмы \ref{alg:topo-indep},
\ref{alg:topo-indep-lookup} и \ref{alg:topo-dep}, учитывающие классы
эквивалентности, существенно уменьшают количество линейных порядков,
необходимых для анализа. Сравнение количества линейных порядков
случайных частично упорядоченных множеств с учётом эквивалентности и
без неё представлено в таблице~\ref{tb:table}. Общее количество
линейных порядков было рассчитано с помощью алгоритма, описанного в
работе~\cite{poset}. Количество линейных порядков с учётом
эквивалентности рассчитано алгоритмом~\ref{alg:topo-indep}.

\begin{table}[ht]
  \captionsetup{justification=raggedright,singlelinecheck=false}
  \caption{Сравнение количества линейных порядков.}
  \centering
  \begin{tabular}{|c|c|c|}
    \hline
    \textbf{Размер множества} & \textbf{
      \begin{tabular}[c]{@{}c@{}}Общее количество \\ линейных порядков
    \end{tabular}} & \textbf{
      \begin{tabular}[c]{@{}c@{}}Количество линейных порядков с
        \\ учётом эквивалентности
    \end{tabular}} \\ \hline
    4                        & 12                                 & 3
    \\ \hline
    9                        & 105                                &
    20
    \\ \hline
    16                       & 2041200                            &
    196
    \\ \hline
    20                       & 1036062752256                      &
    38433
    \\ \hline
    28                       & 6587142659584819200                &
    61881
    \\ \hline
  \end{tabular}\label{tb:table}
\end{table}

При построении таблицы~\ref{tb:table} были выбраны такие множества, в
которых количество рёбер эквивалентности близко к количеству
элементов в частично упорядоченном множестве: $|B^{|}| \sim |P|$.

Некоторые дальнейшие результаты не будут включать алгоритм обхода
топологической сортировкой~\ref{alg:topo}, так как полный обход всех
линейных порядков не может быть произведён за разумное время. Без
эвристики эквивалентности количество линейных порядков в худшем
случае может достигать $n!$. Преимуществом алгоритмов, учитывающих
эквивалентность, является то, что в некоторых пределах
асимптотическая сложность данных алгоритмов может приближаться к
$O(n^2)$ (см.~формулы~\ref{eq:topo-indep-comp-lims},~\ref{eq:ass}).

Обозначим через $p_{D}$ вероятность зависимости между элементами множества $P$:
\begin{equation}
  p_{D} = p_{D}(x, y) = \mathbb{P}[ \{x, y\} \in D ],\ \forall x, y \in P.
  \label{eq:prob_dep}
\end{equation}

Аналогично, через $p_{P}$ обозначим вероятность отношения порядка
между элементами множества $P$:
\begin{equation}
  p_{P} = p_{P}(x, y) = \mathbb{P}[x <_{p} y],\ \forall x, y \in P.
  \label{eq:prob_pos}
\end{equation}

\begin{figure}[ht]
  \centering
  \centerfloat{
    % This file was created with tikzplotlib v0.10.1.post13.
\begin{tikzpicture}

\definecolor{darkgrey176}{RGB}{176,176,176}
\definecolor{green}{RGB}{0,128,0}
\definecolor{lightgrey}{RGB}{211,211,211}
\definecolor{lightgrey204}{RGB}{204,204,204}
\definecolor{orange}{RGB}{255,165,0}

\begin{axis}[
height=0.4\textheight,
legend cell align=right,
legend pos=north west,
legend style={  at={(1,1.03)},  anchor=south east},
tick align=outside,
tick pos=left,
width=0.95\textwidth,
x grid style={darkgrey176},
xlabel={Вероятность зависимости между элементами, \(\displaystyle p_{D}\)},
xmin=-0.0665, xmax=1.0665,
xtick style={color=black},
y grid style={darkgrey176},
ylabel style={text width=10cm, align=center},
ylabel={Количество линейных перестановок
в логарифмическом масштабе, \(\displaystyle log_{10}(\text{шт})\)},
ymin=0, ymax=7.01323097713965,
ytick style={color=black}
]
\draw[draw=black,fill=orange] (axis cs:0.005,0) rectangle (axis cs:0.015,0);
\addlegendimage{ybar,ybar legend,draw=black,fill=orange}
\addlegendentry{Обход с учётом зависимости, Алгоритм~\ref{alg:topo-dep}}

\draw[draw=black,fill=orange] (axis cs:0.105,0) rectangle (axis cs:0.115,2.15836249209525);
\draw[draw=black,fill=orange] (axis cs:0.205,0) rectangle (axis cs:0.215,3.2895889525426);
\draw[draw=black,fill=orange] (axis cs:0.305,0) rectangle (axis cs:0.315,3.74083620705731);
\draw[draw=black,fill=orange] (axis cs:0.405,0) rectangle (axis cs:0.415,4.57801998503221);
\draw[draw=black,fill=orange] (axis cs:0.505,0) rectangle (axis cs:0.515,4.88829746216554);
\draw[draw=black,fill=orange] (axis cs:0.605,0) rectangle (axis cs:0.615,5.21851712477091);
\draw[draw=black,fill=orange] (axis cs:0.705,0) rectangle (axis cs:0.715,5.58294243056767);
\draw[draw=black,fill=orange] (axis cs:0.805,0) rectangle (axis cs:0.815,5.9647525881337);
\draw[draw=black,fill=orange] (axis cs:0.905,0) rectangle (axis cs:0.915,6.28133522204542);
\draw[draw=black,fill=orange] (axis cs:1.005,0) rectangle (axis cs:1.015,6.60437381770785);
\draw[draw=black,fill=white] (axis cs:-0.015,0) rectangle (axis cs:-0.005,4.63578523553365);
\addlegendimage{ybar,ybar legend,draw=black,fill=white}
\addlegendentry{Обход с учётом независимости, Алгоритм~\ref{alg:topo-indep}}

\draw[draw=black,fill=white] (axis cs:0.085,0) rectangle (axis cs:0.095,4.78859255592036);
\draw[draw=black,fill=white] (axis cs:0.185,0) rectangle (axis cs:0.195,4.91948080530445);
\draw[draw=black,fill=white] (axis cs:0.285,0) rectangle (axis cs:0.295,5.118215108501);
\draw[draw=black,fill=white] (axis cs:0.385,0) rectangle (axis cs:0.395,5.44137145388074);
\draw[draw=black,fill=white] (axis cs:0.485,0) rectangle (axis cs:0.495,5.56050560995127);
\draw[draw=black,fill=white] (axis cs:0.585,0) rectangle (axis cs:0.595,5.71155501203113);
\draw[draw=black,fill=white] (axis cs:0.685,0) rectangle (axis cs:0.695,5.88895076288752);
\draw[draw=black,fill=white] (axis cs:0.785,0) rectangle (axis cs:0.795,6.13416817925776);
\draw[draw=black,fill=white] (axis cs:0.885,0) rectangle (axis cs:0.895,6.3165498571675);
\draw[draw=black,fill=white] (axis cs:0.985,0) rectangle (axis cs:0.995,6.60437381770785);
\draw[draw=black,fill=green] (axis cs:-0.005,0) rectangle (axis cs:0.005,0);
\addlegendimage{ybar,ybar legend,draw=black,fill=green}
\addlegendentry{Обход с учётом независимости, Алгоритм~\ref{alg:topo-indep-lookup}}

\draw[draw=black,fill=green] (axis cs:0.095,0) rectangle (axis cs:0.105,2.15836249209525);
\draw[draw=black,fill=green] (axis cs:0.195,0) rectangle (axis cs:0.205,3.2895889525426);
\draw[draw=black,fill=green] (axis cs:0.295,0) rectangle (axis cs:0.305,3.74083620705731);
\draw[draw=black,fill=green] (axis cs:0.395,0) rectangle (axis cs:0.405,4.57801998503221);
\draw[draw=black,fill=green] (axis cs:0.495,0) rectangle (axis cs:0.505,4.88829746216554);
\draw[draw=black,fill=green] (axis cs:0.595,0) rectangle (axis cs:0.605,5.21851712477091);
\draw[draw=black,fill=green] (axis cs:0.695,0) rectangle (axis cs:0.705,5.58294243056767);
\draw[draw=black,fill=green] (axis cs:0.795,0) rectangle (axis cs:0.805,5.9647525881337);
\draw[draw=black,fill=green] (axis cs:0.895,0) rectangle (axis cs:0.905,6.28133522204542);
\draw[draw=black,fill=green] (axis cs:0.995,0) rectangle (axis cs:1.005,6.60437381770785);
\addplot [semithick, red, dash pattern=on 1pt off 3pt on 3pt off 3pt]
table {%
-0.0665 6.60437381770785
1.0665 6.60437381770785
};
\addlegendentry{Обход топологической сортировкой, Алгоритм~\ref{alg:topo}}
\addplot [semithick, lightgrey]
table {%
0 0.17129391428814
0.1 1.90632082152951
0.2 3.14208626146731
0.3 3.98780434183084
0.4 4.5526891703494
0.5 4.94595485475229
0.6 5.26496572446256
0.7 5.55968699567826
0.8 5.86823410629118
0.9 6.22872249419311
1 6.67926759727586
};
\end{axis}

\end{tikzpicture}

  }
  % ./scripts/launch_experiments.py --nds 13 --y-axis=linext
  % --only-max --not-continue --in-log --to-tex
  \caption{График зависимости количества линейных порядков в
    логарифмическом масштабе от вероятности зависимости элементов для
  обходов с учётом зависимости и независимости. $|P| = 13$, $p_{P} = 0.3$.}
  \label{fig:nds-13-linext}
\end{figure}

Из формулы~\ref{eq:supset} следует, что количество линейных порядков
при обходе с учётом $B$ всегда меньше либо равно количеству линейных
порядков при обходе с учётом $B^|$. Кроме того, по построению,
количество линейных порядков всегда равно количеству классов
эквивалентности, что означает минимально возможное количество для
данного частично упорядоченного множества. На
графике~\ref{fig:nds-13-linext} показано, что количество линейных
порядков для одной и той же конфигурации частично упорядоченного
множества меньше при обходе с учётом зависимостей по сравнению с
обходом с учётом независимостей. Красной линией обозначено
максимально возможное количество линейных порядков для данной
конфигурации частично упорядоченного множества.

\begin{figure}[ht]
  \centering
  \centerfloat{
    % This file was created with tikzplotlib v0.10.1.post13.
\begin{tikzpicture}

  \definecolor{darkgrey176}{RGB}{176,176,176}
  \definecolor{green}{RGB}{0,128,0}
  \definecolor{lightgrey}{RGB}{211,211,211}
  \definecolor{lightgrey204}{RGB}{204,204,204}
  \definecolor{orange}{RGB}{255,165,0}

  \begin{axis}[
      height=0.4\textheight,
      legend cell align=right,
      legend pos=north west,
      legend style={  at={(1,1.03)},  anchor=south east},
      tick align=outside,
      tick pos=left,
      width=0.95\textwidth,
      x grid style={darkgrey176},
      xlabel={Вероятность зависимости между элементами,
      \(\displaystyle p_{D}\)},
      xmin=-0.0665, xmax=1.0665,
      xtick style={color=black},
      y grid style={darkgrey176},
      ylabel style={text width=10cm, align=center},
      ylabel={Время
      в логарифмическом масштабе, \(\displaystyle log_{10}(\text{мс} * 100)\)},
      ymin=0, ymax=6.45493824874219,
      ytick style={color=black}
    ]
    \draw[draw=black,fill=orange] (axis cs:0.005,0) rectangle (axis
    cs:0.015,0.672698063024046);
    \addlegendimage{ybar,ybar legend,draw=black,fill=orange}
    \addlegendentry{Обход с учётом зависимости, Алгоритм~\ref{alg:topo-dep}}

    \draw[draw=black,fill=orange] (axis cs:0.105,0) rectangle (axis
    cs:0.115,2.200896988085);
    \draw[draw=black,fill=orange] (axis cs:0.205,0) rectangle (axis
    cs:0.215,3.10438808634469);
    \draw[draw=black,fill=orange] (axis cs:0.305,0) rectangle (axis
    cs:0.315,3.30237422607558);
    \draw[draw=black,fill=orange] (axis cs:0.405,0) rectangle (axis
    cs:0.415,4.16833245562574);
    \draw[draw=black,fill=orange] (axis cs:0.505,0) rectangle (axis
    cs:0.515,4.46480837004647);
    \draw[draw=black,fill=orange] (axis cs:0.605,0) rectangle (axis
    cs:0.615,4.75409670845534);
    \draw[draw=black,fill=orange] (axis cs:0.705,0) rectangle (axis
    cs:0.715,5.12614440017105);
    \draw[draw=black,fill=orange] (axis cs:0.805,0) rectangle (axis
    cs:0.815,5.48329053467323);
    \draw[draw=black,fill=orange] (axis cs:0.905,0) rectangle (axis
    cs:0.915,5.80480570279117);
    \draw[draw=black,fill=orange] (axis cs:1.005,0) rectangle (axis
    cs:1.015,6.07932598675282);
    \draw[draw=black,fill=white] (axis cs:-0.015,0) rectangle (axis
    cs:-0.005,3.91260621327443);
    \addlegendimage{ybar,ybar legend,draw=black,fill=white}
    \addlegendentry{Обход с учётом независимости, Алгоритм~\ref{alg:topo-indep}}

    \draw[draw=black,fill=white] (axis cs:0.085,0) rectangle (axis
    cs:0.095,4.03298119309737);
    \draw[draw=black,fill=white] (axis cs:0.185,0) rectangle (axis
    cs:0.195,4.13142296322114);
    \draw[draw=black,fill=white] (axis cs:0.285,0) rectangle (axis
    cs:0.295,4.27295494484817);
    \draw[draw=black,fill=white] (axis cs:0.385,0) rectangle (axis
    cs:0.395,4.47818686310984);
    \draw[draw=black,fill=white] (axis cs:0.485,0) rectangle (axis
    cs:0.495,4.60648992516164);
    \draw[draw=black,fill=white] (axis cs:0.585,0) rectangle (axis
    cs:0.595,4.72960453322533);
    \draw[draw=black,fill=white] (axis cs:0.685,0) rectangle (axis
    cs:0.695,4.87760686083614);
    \draw[draw=black,fill=white] (axis cs:0.785,0) rectangle (axis
    cs:0.795,5.09484132652886);
    \draw[draw=black,fill=white] (axis cs:0.885,0) rectangle (axis
    cs:0.895,5.32629693656581);
    \draw[draw=black,fill=white] (axis cs:0.985,0) rectangle (axis
    cs:0.995,5.52179034348834);
    \draw[draw=black,fill=green] (axis cs:-0.005,0) rectangle (axis
    cs:0.005,2.60540848858583);
    \addlegendimage{ybar,ybar legend,draw=black,fill=green}
    \addlegendentry{Обход с учётом независимости,
    Алгоритм~\ref{alg:topo-indep-lookup}}

    \draw[draw=black,fill=green] (axis cs:0.095,0) rectangle (axis
    cs:0.105,3.17469347914945);
    \draw[draw=black,fill=green] (axis cs:0.195,0) rectangle (axis
    cs:0.205,3.39158253515756);
    \draw[draw=black,fill=green] (axis cs:0.295,0) rectangle (axis
    cs:0.305,3.88808902988961);
    \draw[draw=black,fill=green] (axis cs:0.395,0) rectangle (axis
    cs:0.405,4.20617259446969);
    \draw[draw=black,fill=green] (axis cs:0.495,0) rectangle (axis
    cs:0.505,4.40737900582878);
    \draw[draw=black,fill=green] (axis cs:0.595,0) rectangle (axis
    cs:0.605,4.64942054065054);
    \draw[draw=black,fill=green] (axis cs:0.695,0) rectangle (axis
    cs:0.705,4.89481322245816);
    \draw[draw=black,fill=green] (axis cs:0.795,0) rectangle (axis
    cs:0.805,5.13628288216488);
    \draw[draw=black,fill=green] (axis cs:0.895,0) rectangle (axis
    cs:0.905,5.3313847760984);
    \draw[draw=black,fill=green] (axis cs:0.995,0) rectangle (axis
    cs:1.005,5.4919762757976);
    \addplot [semithick, lightgrey]
    table {%
      0 0.813095484880913
      0.1 2.01027690596296
      0.2 2.91956692726929
      0.3 3.59800273858807
      0.4 4.10262152970745
      0.5 4.49046049041558
      0.6 4.8143432165528
      0.7 5.11023892816826
      0.8 5.40990325136328
      0.9 5.74509181223919
      1 6.14756023689733
    };
  \end{axis}

\end{tikzpicture}

  }
  % ./scripts/launch_experiments.py --nds 13 --y-axis=time --only-max
  % --not-continue --in-log --to-tex
  \caption{График зависимости времени построения линейных порядков в
    логарифмическом масштабе от вероятности зависимости элементов для
  обходов с учётом зависимости и независимости. $|P| = 13$, $p_{P} = 0.3$.}
  \label{fig:nds-13-time}
\end{figure}

Как видно из графика~\ref{fig:nds-13-time}, выбор
оптимального
алгоритма построения линейных порядков зависит от количества
зависимых пар элементов в частично упорядоченном множестве. При
значениях параметра $p_{D} \lesssim 0.5$ было показано, что учёт
зависимостей обеспечивает существенное преимущество в
производительности по сравнению с учётом независимости. Поскольку оба
подхода (см. алгоритмы~\ref{alg:topo-dep}
и~\ref{alg:topo-indep-lookup}) реализуют построение классов
эквивалентных линейных порядков, выбор конкретного алгоритма
определяется структурой исходного частично упорядоченного множества.

\subsection{Применение алгоритма построения линейных порядков с
учётом зависимости к тестовым сценариям TSOTool}

Преимущество обхода частично упорядоченного множества с учётом
зависимостей (см. алгоритм~\ref{alg:topo-dep}) наглядно проявляется на
тестовых сценариях инструмента TSOTool (см. раздел~\ref{sec:manovit}).

Частично упорядоченное множество определяется моделью памяти. Если
результат операции $b$ не может наблюдаться до выполнения операции
$a$, то $a < b$. Пара ``адрес/значение'' задаёт отношение зависимости
между операциями. Условие уникальности пар ``адрес/значение'' при
верификации с использованием алгоритма~\ref{alg:topo-dep} однозначно
определяет порядок между операциями в момент выбора направления (либо
$P^{\rightarrow}$, либо $P^{\leftarrow}$).

\begin{figure}[ht]
  \centering
  \centerfloat{
    % This file was created with tikzplotlib v0.10.1.post13.
\begin{tikzpicture}

  \definecolor{darkgrey176}{RGB}{176,176,176}
  \definecolor{lightgrey204}{RGB}{204,204,204}

  \begin{axis}[
      height=0.4\textheight,
      legend cell align=right,
      legend pos=north west,
      legend style={  at={(1,1.03)},  anchor=south east},
      tick align=outside,
      tick pos=left,
      width=0.95\textwidth,
      x grid style={darkgrey176},
      xlabel={\(\displaystyle |P|, шт\)},
      xmin=6800, xmax=33200,
      xtick style={color=black},
      y grid style={darkgrey176},
      ylabel style={text width=10cm, align=center},
      ylabel={Время, мс},
      ymin=-275.5095, ymax=5785.6995,
      ytick style={color=black}
    ]
    \path [draw=black, semithick]
    (axis cs:8000,0)
    --(axis cs:8000,261.365);

    \addplot [semithick, black, mark=*, mark size=3, mark
    options={solid}, only marks]
    table {%
      8000 261.365
    };
    \addplot [semithick, red]
    table {%
      8000 0
      8000 0
    };
    \path [draw=black, semithick]
    (axis cs:12000,0)
    --(axis cs:12000,638.937);

    \addplot [semithick, black, mark=*, mark size=3, mark
    options={solid}, only marks]
    table {%
      12000 638.937
    };
    \addplot [semithick, red]
    table {%
      12000 0
      12000 0
    };
    \path [draw=black, semithick]
    (axis cs:16000,0)
    --(axis cs:16000,1137.16);

    \addplot [semithick, black, mark=*, mark size=3, mark
    options={solid}, only marks]
    table {%
      16000 1137.16
    };
    \addplot [semithick, red]
    table {%
      16000 0
      16000 0
    };
    \path [draw=black, semithick]
    (axis cs:20000,0)
    --(axis cs:20000,1866.19);

    \addplot [semithick, black, mark=*, mark size=3, mark
    options={solid}, only marks]
    table {%
      20000 1866.19
    };
    \addplot [semithick, red]
    table {%
      20000 0
      20000 0
    };
    \path [draw=black, semithick]
    (axis cs:24000,0)
    --(axis cs:24000,2700.01);

    \addplot [semithick, black, mark=*, mark size=3, mark
    options={solid}, only marks]
    table {%
      24000 2700.01
    };
    \addplot [semithick, red]
    table {%
      24000 0
      24000 0
    };
    \path [draw=black, semithick]
    (axis cs:28000,0)
    --(axis cs:28000,4113.57);

    \addplot [semithick, black, mark=*, mark size=3, mark
    options={solid}, only marks]
    table {%
      28000 4113.57
    };
    \addplot [semithick, red]
    table {%
      28000 0
      28000 0
    };
    \path [draw=black, semithick]
    (axis cs:32000,0)
    --(axis cs:32000,5510.19);

    \addplot [semithick, black, mark=*, mark size=3, mark
    options={solid}, only marks]
    table {%
      32000 5510.19
    };
    \addplot [semithick, red]
    table {%
      32000 0
      32000 0
    };
  \end{axis}

\end{tikzpicture}

  }
  \caption{График времени построения одного линейного порядка в
  зависимости от количества элементов множества. $p_{P} = 0.001$.}
  \label{fig:manovit-exp}
  % ./build/examples/example --algorithm dep --nodes 28000
  % --poset-edge-prob 0.001 --dependency-edge-prob 0.3 --seed 0
  % --predefined 3 --linexts-act print-one
  % ./scripts/manovit-exp.py --nds 8000 12000 16000 20000 24000 28000
  % 32000 --pepr 0.001 --experiments 1 --only-max --not-continue
\end{figure}

Промоделируем эксперимент с графом, получаемым инструментом TSOTool.
Предположим, что половина всех операций с памятью — это операции
чтения, а другая половина — операции записи. Во время верификации
цикл между операциями $M[x] == v$ и $M[x] := v$ возможен только в
случае некорректного поведения тестируемой системы.
Алгоритм~\ref{alg:topo-dep} допускает циклы, если возможны другие
порядки чтения/записи, однако в контексте построения псевдослучайного
теста такие порядки невозможны. Для моделирования отсутствия циклов
используем сильно разреженное частично упорядоченное множество с
параметром $p_{P} = 0.001$. Уникальность пар ``адрес/значение''
моделируется наличием только одного ребра между операцией чтения и
операцией записи.

На графике~\ref{fig:manovit-exp} представлены результаты построения
одного порядка в зависимости от количества элементов множества. Как
видно, обход с учётом зависимостей может применяться для проверки
нескольких тысяч операций с памятью.

\clearpage
